\section{Hai thứ lắp thêm lúc giải mã}
\label{sec:ch05-giai-ma}

\index{beam search!bài toán nó giải}%
Phương trình~\eqref{eq:ch05-phan-ra} cho xác suất của một
\tn{câu đích}{target sentence} cho trước, khi đã có
\tn{câu nguồn}{source sentence} và \tn{vector ngữ cảnh}{context vector} rút ra
từ nó. Dịch thì cần chiều ngược lại: tìm câu đích có xác suất cao nhất,

\begin{equation}
  \hat{T} = \argmax_{T} p(T \mid S),
  \label{eq:ch05-argmax}
\end{equation}

và phép \(\argmax\) ấy chạy trên mọi chuỗi token thuộc mọi độ dài. Không tính
được. Lấy từ có xác suất cao nhất ở mỗi bước, tức giải mã tham lam, thì tính
được nhưng không giải \eqref{eq:ch05-argmax}: một từ rẻ ở bước này có thể ép mọi
bước sau vào một nhánh đắt.

\subsection*{Beam search}

\index{beam search!thuật toán}%
Beam search là chỗ ở giữa. Giữ \(B\) \tn{giả thuyết}{hypothesis} tốt nhất thay
vì một. Bài mô tả nó trong mục 3.2, và bản giả mã dưới đây chép theo mô tả ấy;
bản chạy được là \code{beam\_decode} trong repo companion.

\begin{algorithm}[htbp]
  \caption{Beam search, giả mã theo mô tả ở mục 3.2 của bài báo}
  \label{alg:ch05-beam}
  \begin{algorithmic}[1]
    \State \(\mathcal{L} \gets \{(\varepsilon,\ 0)\}\),
           \(\mathcal{F} \gets \emptyset\)
    \While{\(\mathcal{L} \neq \emptyset\)}
      \State \(\mathcal{C} \gets \{(\vect{u} \cdot w,\ s + \log p(w \mid \vect{u}))
              : (\vect{u}, s) \in \mathcal{L},\ w \in V\}\)
      \State giữ \(B\) phần tử điểm cao nhất của \(\mathcal{C}\)
      \State chuyển phần tử kết thúc bằng \(\langle\)EOS\(\rangle\) sang
             \(\mathcal{F}\); số còn lại thành \(\mathcal{L}\)
      \State bỏ khỏi \(\mathcal{L}\) mọi giả thuyết có điểm thấp hơn
             \(\max \mathcal{F}\)
    \EndWhile
    \State \Return phần tử điểm cao nhất của \(\mathcal{F}\)
  \end{algorithmic}
\end{algorithm}

Dòng 6 là dòng tôi phải sửa lại sau khi đo, và nó đáng kể ra vì cái sai của nó
không hiện ra thành lỗi. Bản đầu tôi viết dừng ngay khi \(\mathcal{F}\) đủ
\(B\) phần tử, rồi vứt hết \(\mathcal{L}\). Nghe hợp lý, và sai: một giả thuyết
dài hơn còn đang sống có thể đang dẫn điểm so với cả \(B\) giả thuyết vừa kết
thúc.

\begin{measured}{luật dừng sai, đo lại trên đúng cấu hình của tag \code{ch05}}
\begin{minted}{text}
stop rule  beam  exact    short    long     mean output length
greedy     -     0.8533   1.0000   0.7105   13.883
broken     2     0.8400   0.9527   0.7303   13.827
broken     5     0.8767   1.0000   0.7566   13.863
broken     12    0.8800   1.0000   0.7632   13.867
fixed      2     0.8733   1.0000   0.7500   13.887
fixed      5     0.8833   1.0000   0.7697   13.877
\end{minted}

Bản \code{broken} là \code{beam\_decode} với dòng 6 thay bằng
\enquote{dừng khi \(\mathcal{F}\) đủ \(B\) phần tử}, tức đúng phép sửa ngược mà
bài tập~3 phần \enquote{Để tay vào} yêu cầu bạn dựng lại. Nó không nằm trong
tag; hàm đầy đủ nằm trong research note của chương.
\end{measured}

Ở beam 2, luật sai cho 0.8400 trong khi giải mã tham lam cho 0.8533. Beam rộng
hơn mà lại thua beam bằng một, và chỗ mất nằm gọn ở câu \emph{ngắn}: 1.0000 tụt
xuống 0.9527, còn câu dài thì lại nhích lên. Đúng chiều mà cơ chế dự đoán, vì
câu ngắn là câu có giả thuyết kết thúc sớm, nên chúng là nhóm bị luật dừng cắt
ngang.

Bản ở dòng 6 thì cắt tỉa theo một bất đẳng thức chứ không theo một số đếm. Log
xác suất chỉ có giảm, nên một giả thuyết còn sống mà đã thua giả thuyết hoàn tất
tốt nhất thì không bao giờ đuổi kịp được nữa. Cắt nó đi là phép cắt chính xác,
không phải phép xấp xỉ, và nó cho tìm kiếm dừng sớm mà không mất gì: cùng mô
hình, beam 2 lên 0.8733 và beam 5 lên 0.8833.

\begin{measured}{độ chính xác theo bề rộng beam, một mô hình}
\begin{minted}{text}
beam  exact    short    long     mean output length
1     0.8533   1.0000   0.7105   13.883
2     0.8733   1.0000   0.7500   13.887
5     0.8833   1.0000   0.7697   13.877
12    0.8833   1.0000   0.7697   13.877
\end{minted}

Mô hình seed 0, câu nguồn đảo, 300 câu test. Độ dài câu đích tham chiếu trung
bình là 13.877.

Đo bằng \code{experiments/ch05\_search.py} tại tag \code{ch05}.
\end{measured}

\index{beam search!lợi ích bão hòa rất sớm}%
Hình dạng khớp với điều bài nói: hệ chạy tốt ngay ở beam 1, và beam 2 đã lấy
được gần hết phần mà beam search mang lại. Nguyên văn:
\enquote{a beam of size 2 provides most of the benefits of beam search}. Ở đây
beam 2 lấy được 0.0200 trên tổng 0.0300 mà beam mang lại, tức hai phần ba, và
beam 12 không hơn beam 5 một chữ số nào.

\subsection*{Một cách đọc sai, và cách số liệu bác bỏ nó}

\index{beam search!thiên lệch độ dài}%
Lúc beam 2 còn thua tham lam, cách đọc đầu tiên của tôi là thiên lệch độ dài:
điểm của một giả thuyết là tổng các log xác suất, toàn số âm, nên câu càng dài
điểm càng thấp và beam search sẽ chuộng câu ngắn. Đây là một hiện tượng có thật
và có tên, nên nó dễ tin.

Nó sai ở đây, và cột cuối của bảng trên là chỗ bác bỏ. Ở luật dừng sai, độ dài
output trung bình là 13.827 tại beam 2, so với 13.877 của câu tham chiếu. Lệch
0.05 token, không phải một xu hướng chuộng câu ngắn. Tôi có thử thêm phép chuẩn
hoá theo độ dài, tức chia điểm cho số token trước khi so, và cả ba cột 0.8400,
0.9527 và 0.7303 đứng nguyên không đổi một chữ số nào; thứ duy nhất nhúc nhích
là độ dài trung bình, 13.827 thành 13.840. Chuẩn hoá độ dài là bản vá cho một
lỗi corpus này không có, và tôi đã suýt áp nó vì cái tên nghe đúng.

Còn một chuyện về beam search mà tôi tin sai theo chiều khác, và ghi lại vì nó
là một cái bẫy sạch sẽ. Beam rộng sinh ra tập ứng viên bao tập của beam hẹp, nên
có vẻ hiển nhiên là nó phải trả về giả thuyết có log xác suất không thấp hơn.
Không đúng: mỗi bước chỉ giữ \(B\) ứng viên tốt nhất, và tiền tố mà beam hẹp giữ
lại có thể xếp dưới hạng \(B\) trong tập ứng viên lớn hơn của beam rộng, nên bị
loại mất. Tôi viết hẳn một test khẳng định tính đơn điệu ấy trước khi phát hiện
ra nó sai; test bây giờ khẳng định đúng thứ nó bảo vệ, và phản ví dụ nằm trong
research note. Nói cách khác, một cột beam đi xuống ở đâu đó không phải bằng
chứng có bug.

\subsection*{Ensemble}

\index{ensemble!cách gộp}%
Thứ thứ hai lắp thêm lúc giải mã là ensemble: huấn luyện vài mô hình rồi cho
chúng cùng quyết định. Bài dùng năm LSTM khác nhau ở khởi tạo ngẫu nhiên và ở
thứ tự ngẫu nhiên của các minibatch (\enquote{that differ in their random
initializations and in the random order of minibatches}), và cả hai thứ ấy ở
đây đều do một seed điều khiển.

Cách gộp là chỗ đáng nói rõ. Ensemble không phải một cuộc bỏ phiếu chọn câu cuối
cùng. Ở mỗi bước giải mã, log xác suất của các mô hình được lấy trung bình
\emph{trước} khi beam search xếp hạng, nên mọi thành viên đều góp phần nặn ra
từng bước một chứ không phải chọn giữa vài bản dịch làm sẵn.

\begin{measured}{ba mô hình chỉ khác nhau ở seed, và ensemble của chúng}
\begin{minted}{text}
one model at a time, beam 2:
seed  exact    short    long
0     0.8733   1.0000   0.7500
1     0.6767   0.9797   0.3816
2     0.8700   1.0000   0.7434

ensemble size, beam 2:
models  exact    short    long
1       0.8733   1.0000   0.7500
2       0.9467   1.0000   0.8947
3       0.9867   1.0000   0.9737
\end{minted}

Đo bằng \code{experiments/ch05\_search.py} tại tag \code{ch05}.
\end{measured}

\index{ensemble!so với beam search}%
Bảng trên có hai chuyện, và chuyện thứ nhất không phải chuyện về ensemble.

Ba mô hình ấy khác nhau đúng một con số seed. Chúng cho 0.8733, 0.6767 và
0.8700. Gần hai mươi điểm tản, trên cùng dữ liệu, cùng kiến trúc, cùng số bước
cập nhật. Đây là bối cảnh mà mọi con số khác trong chương phải được đọc cùng, và
là lý do mục~\ref{sec:ch05-dao-cau-nguon} phải chạy ba seed rồi so theo cặp.

Chuyện thứ hai: đặt hai phép cải thiện cạnh nhau. Nới beam từ 1 lên 12 được
0.0300. Thêm hai mô hình được 0.1134, gấp gần bốn lần, và trên câu dài thì từ
0.7500 lên 0.9737. Ensemble ba mô hình cũng cao hơn mô hình đơn \emph{tốt nhất}
trong ba, nên đây không phải chuyện gộp lại để đỡ rút phải mô hình xấu.

Bài cho cùng thứ tự độ lớn ở quy mô của nó. Một LSTM đảo với beam 12 được BLEU
30.59; ensemble năm LSTM với beam 12 được 34.81. Còn nới beam từ 1 lên 12 trên
cùng ensemble năm mô hình chỉ đi từ 33.00 lên 34.81. Chú thích bảng 1 của bài
rút ra kết luận thực dụng từ đúng chỗ đó: một ensemble năm LSTM chạy beam 2 rẻ
hơn một LSTM đơn chạy beam 12, và nó cho điểm cao hơn.

Cả hai thứ trong mục này đều không đụng tới cái mà
mục~\ref{sec:ch05-bop-vector} đo. Chúng lấy được nhiều hơn từ những vector ngữ
cảnh đã có, chứ không nới rộng vector nào cả.
